\unsection{ВСТУП}

\subsection*{Актуальність, мета та постановка задачі}

Повнотекстовий пошук є типовою задачею, де паралелізм потрібен на двох рівнях одночасно.
По-перше, побудова інвертованого індексу над великою колекцією -- це обчислювально важка
пакетна робота: кожен байт тексту треба прочитати, розбити на слова й врахувати у
спільній структурі даних. По-друге, готовий індекс обслуговує багатьох клієнтів одночасно, і
сервер має обробляти їхні запити паралельно, не блокуючи один одного. Основне завдання
паралельної обробки -- скорочення часу виконання програми [1]; у розробленій системі
побудова індексу корпусу з 16\,793 документів на одному потоці триває 10,1~с, а на десяти --
2,1~с (підрозділ~3.4). Третій чинник -- інтеграція великої мовної моделі для розширення
запитів: модель на CPU відповідає за сотні мілісекунд і стає послідовною ділянкою конвеєра,
яку доводиться ізолювати від решти системи.

\textit{Метою роботи} є дослідження та практична реалізація паралельних обчислень на
прикладі пошукової системи з інтеграцією LLM: паралельної побудови інвертованого індексу за
допомогою власного пулу потоків, потокобезпечного доступу до спільної структури індексу та
одночасної обробки запитів від багатьох клієнтів.

\textit{Об'єкт дослідження} -- процеси паралельної побудови інвертованого індексу та
паралельної обробки пошукових запитів. \textit{Предмет дослідження} -- механізми
синхронізації потоків у пошуковій системі (блокувальні та неблокуючі способи злиття
результатів у спільну структуру, безблокувальне читання незмінного знімка індексу, обмеження
паралельних звернень до моделі) та їхній вплив на час побудови, пропускну здатність і
використання пам'яті. \textit{Методи дослідження} -- повторні вимірювання з випадковим
чергуванням конфігурацій, аналіз прискорення за законом Амдала та метрикою Карпа -- Флатта
(підрозділ~3.1), розклад часу за фазами, порівняльний експеримент «до/після» оптимізації,
модульне, стресове та наскрізне тестування.

Відповідно до завдання система має задовольняти такі вимоги:
\begin{itemize}
\item бекенд компільованою мовою (обрано Java 21), два різні клієнти -- консольний і веб;
\item запуск однією командою \code{docker compose up}: сервіс пошукового сервера і сервіс
Ollama, ендпоінт \code{GET /api/health};
\item самостійно реалізовані пул потоків і структура інвертованого індексу з паралельним
доступом, кількість потоків побудови задається параметром \code{threads};
\item фіксований REST-контракт: \code{GET /api/health}, \code{POST /api/index},
\code{POST /api/search} (режим \code{keyword}), \code{GET /api/stats}, формат JSON;
\item розширення запиту ключовими словами та синонімами через Ollama \code{/api/generate};
\item WebSocket \code{/ws/search}, що передає етапи \code{query\_expansion},
\code{keyword\_search}, \code{done} окремими повідомленнями з полями \code{stage} і
\code{data};
\item текстовий датасет англійською мовою обсягом не менше 1~ГБ і не менше 10\,000 документів;
\item графіки часу побудови індексу від кількості потоків і пропускної здатності від
кількості клієнтів, аналіз вузьких місць, UML-діаграми, інструкція зі збирання, історія
комітів у стилі Conventional Commits.
\end{itemize}

\subsection*{Обрані технології}

\textit{Java 21 без сторонніх бібліотек.} Мова надає модель пам'яті з чітко визначеним
відношенням happens-before [15], атомарні змінні з операцією CAS (\code{AtomicReference},
\code{AtomicBoolean}) та примітиви \code{ReentrantLock}/\code{Condition} [4] -- саме ті засоби,
порівняння яких є предметом роботи. Усі компоненти, що зазвичай беруть із бібліотек, --
пул потоків, інвертований індекс, HTTP/1.1-сервер, протокол WebSocket, розбір JSON,
навантажувальний клієнт і навіть тестовий фреймворк -- реалізовано в проєкті. Це
спрощує перевірку відсутності станів гонки (весь код синхронізації видимий) і дозволяє
збирати проєкт лише засобами JDK (\code{javac} і \code{jar}) без доступу до репозиторіїв
залежностей.

\textit{Docker і Docker Compose} [17] забезпечують відтворюване розгортання: сервер, Ollama та
два одноразові сервіси ініціалізації (розпакування корпусу і завантаження моделі)
піднімаються однією командою.

\textit{Ollama і модель qwen2.5:1.5b} [16, 21]. Модель обрано як компроміс між якістю
синонімів і затримкою на CPU: вона дає змістовні розширення (табл.~3.10) з медіанною
затримкою 0,64~с (підрозділ~3.9); більші моделі тієї самої родини потребують пропорційно
більше обчислень на кожен згенерований токен і на CPU відповідали б повільніше.

\textit{Датасет} -- архів книжок Project Gutenberg [20] обсягом 1,70~ГБ (2\,809 файлів). Оскільки
книжок менше за потрібні 10\,000 документів, кожну книжку розрізано по межах абзаців на
документи приблизно по 100~КБ (підрозділ~2.1), що дало 16\,793 документи обсягом 1,62~ГБ.
У силабусі в одному місці вказано обсяг «не менше 2~ГБ», а у вимогах до датасету -- «не менше
1~ГБ»; робота виконує вимогу розділу «Вимоги до датасету».

\textit{Допоміжні засоби}: PlantUML для UML-діаграм, Python з matplotlib для побудови
графіків із CSV-результатів, Git і GitHub для керування версіями. Повідомлення комітів
оформлено за специфікацією Conventional Commits [18], код -- за Google Java Style [19].
